\apendice{Plan de Proyecto Software}

\section{Introducción}
Este documento define cómo se va a gestionar y ejecutar el desarrollo de un producto o sistema informático. Establece el alcance, los objetivos, el cronograma, los recursos, el presupuesto, los riesgos y los criterios de calidad, así como la organización del equipo y los mecanismos de seguimiento y control. En esencia, convierte una idea de software en un itinerario gestionable, con responsabilidades claras y procedimientos para controlar el progreso y los cambios.

Su importancia radica en que reduce la incertidumbre y aumenta la probabilidad de éxito académico y técnico. Permite justificar decisiones, anticipar problemas, y demostrar competencias en gestión de proyectos, ingeniería de requisitos, calidad y comunicación técnica.

En este documento se va a incluir la planificación temporal de las fases y tareas del proyecto. Para cada una se debe establecer una fecha de inicio y se debe estimar la fecha final. Para realizar esto hay que tener en cuenta el peso de cada tarea, los requisitos previos y las dependencias.

También, se va a valorar la viabilidad del proyecto. En concreto, la viabilidad económica y la legal. En la viabilidad económica se estiman los costes y los beneficios del proyecto. Mientras que en la viabilidad legal analizan las leyes y licencias que afecten al desarrollo.  

\section{Planificación temporal}
Para la planificación se utiliza una metodología Scrum, aplicando iteraciones (\eng{sprints}). Al finalizar cada \eng{sprint} se realiza una reunión para la revisión del incremento y la planificación del siguiente. En la planificación se seleccionan y valoran las tareas (\eng{issues}) a realizar en el siguiente \eng{sprint}. Estas tareas se han valorado usando los \eng{story points} de ZenHub. A las tareas que se subdividen en otras se les ha asignado solo puntos extra respecto a sus subtareas, los \eng{story points} no se han considerado acumulativos.

Para simplificar la asignación de los \eng{story points} se les ha asignado una estimación temporal recogida en la tabla~\ref{tabla:storypoints}:

\begin{table}[h]
	\centering
	\begin{tabularx}{\linewidth}{
    		>{\centering\arraybackslash}p{0.46\linewidth}
    		>{\centering\arraybackslash}p{0.46\linewidth}}
		\toprule
		\textbf{\eng{Story Points}}    & \textbf{Tiempo estimado (horas)}\\
		\midrule
		1 & 0,5  \\
		2 & 1  \\
		3 & 2  \\
		5 & 4  \\
		8 & 6  \\
		13 & 9  \\
		21 & 15  \\
		40 & más de 20 \\
		\bottomrule
	\end{tabularx}
	\caption{Tiempo estimado de los \eng{story points}.}
	\label{tabla:storypoints}
\end{table}


\subsection{\eng{Sprint} 0 (11/09/2025 -- 18/09/2025)}
Este \eng{sprint} marcó el comienzo del proyecto. Los objetivos fueron realizar una primera toma de contacto con los \eng{RAG} y los \eng{LLM}, investigando sobre qué son, cómo funcionan y qué ventajas aportan los \eng{RAG}. También, se planteó empezar a escribir los conceptos teóricos relevantes y crear el repositorio \myhref{https://github.com/lbr1014/TFG_RAG.git}{GitHub} para el proyecto. Como el \eng{sprint} se diseñó antes de la creación del repositorio, no se creó el \eng{milestone} asociado como se ha hecho en los \eng{sprints} posteriores. 
Para llevar a cabo estos objetivos se plantearon las siguientes tareas:
\begin{itemize}
\tightlist
	\item Crear el repositorio en GitHub.
	\item Investigar sobre los \eng{RAG} y \eng{LLM}.
	\begin{itemize}
	\tightlist
		\item Leer el capítulo 1 del libro: \eng{A simple guide to retrieval augmented generation}~\cite{book:RAG_SimpleGuide}.
		\item Leer el capítulo 2 del libro: \eng{A simple guide to retrieval augmented generation}~\cite{book:RAG_SimpleGuide}
	\end{itemize}
	\item Escribir conceptos teóricos sobre lo investigado.
\end{itemize}
Para realizar estas tareas se estimaron 6 horas de trabajo pero en realidad fueron 9. Al finalizar el \eng{sprint} se completaron todas las tareas.

\subsection{\eng{Sprint} 1 (18/09/2025 -- 25/09/2025)}
Los objetivos del \myhref{https://github.com/lbr1014/TFG_RAG/milestone/1?closed=1}{Sprint 1} fueron investigar sobre los \eng{LLM} y los \eng{RAG} y escribir la información más relevante como conceptos teóricos en la memoria. Comenzar a escribir los objetivos del proyecto y las técnicas y herramientas definidas hasta el momento, como la metodología Scrum o el repositorio. 
Para llevar a cabo estos objetivos se plantearon las siguientes tareas:
\begin{itemize}
\tightlist
	\item Escribir objetivos del proyecto y técnicas y herramientas.
	\item Escribir conceptos teóricos.
	\begin{itemize}
	\tightlist
		\item Investigar sobre los \eng{RAG} y los \eng{LLM}.
		\begin{itemize}
		\tightlist
			\item Leer el capítulo 3 del libro: \eng{A simple guide to retrieval augmented generation}~\cite{book:RAG_SimpleGuide}.
			\item Leer el capítulo 4 del libro: \eng{A simple guide to retrieval augmented generation}~\cite{book:RAG_SimpleGuide}.
		\end{itemize}
	\end{itemize}
\end{itemize}

Para realizar estas tareas se estimaron 8,5 horas de trabajo pero en realidad fueron 10. Al finalizar el \eng{sprint} se completaron todas las tareas.

El \eng{burndown} de este \eng{sprint} se ilustra en la figura~\ref{fig:Burndown/Sprint1}.
\imagenflotante{Burndown/Sprint1}{\eng{Burndown chart} del \eng{sprint} 1.}

\subsection{\eng{Sprint} 2 (25/09/2025 -- 02/10/2025)}
Los objetivos del \myhref{https://github.com/lbr1014/TFG_RAG/milestone/2?closed=1}{Sprint 2} fueron seguir investigando sobre los \eng{LLM} y \eng{RAG} y explicar la información relevante en los conceptos teóricos. Investigar sobre el \eng{Web Scraping} y explicar en los anexos el plan de proyecto, en concreto la planificación temporal y los \eng{sprints}.
Para llevar a cabo estos objetivos se plantearon las siguientes tareas:
\begin{itemize}
\tightlist
	\item Añadir las correcciones a la memoria.
	\item Completar el apartado \eng{pipeline} de generación del apartado conceptos teóricos.
	\item Investigar sobre \eng{Web Scraping}.
	\begin{itemize}
	\tightlist
		\item Investigar sobre HTML.
		\item Investigar sobre CSS.
	\end{itemize}
	\item Explicar el plan de proyectos en los anexos.
\end{itemize}

Para realizar estas tareas se estimaron 10 horas de trabajo pero en realidad fueron 12. Al finalizar el \eng{sprint} se completaron todas las tareas.

El \eng{burndown} de este \eng{sprint} se ilustra en la figura~\ref{fig:Burndown/Sprint2}.
\imagenflotante{Burndown/Sprint2}{\eng{Burndown chart} del \eng{sprint} 2.}

\subsection{\eng{Sprint} 3 (2/10/2025 -- 09/10/2025)}
Los objetivos del \myhref{https://github.com/lbr1014/TFG_RAG/milestone/3?closed=1}{Sprint 3} fueron realizar un primer prototipo de \eng{Web Scraping}, tanto en \eng{Selenium} como en \eng{Playwright}. Incluir en la memoria las correcciones y la información relevante.
Para llevar a cabo estos objetivos se plantearon las siguientes tareas:
\begin{itemize}
\tightlist
	\item Añadir las correcciones a la memoria.
	\item Realización del comienzo de la aplicación de \eng{web scraping}.
	\begin{itemize}
	\tightlist
		\item Investigar \eng{web scraping} con Selenium.
		\item Investigar \eng{web scraping} con Playwright.
	\end{itemize}
	\item Escribir conceptos teóricos en la memoria.
	\begin{itemize}
	\tightlist
		\item Leer el capítulo 5 del libro: \eng{A simple guide to retrieval augmented generation}~\cite{book:RAG_SimpleGuide}.
	\end{itemize}
\end{itemize}

Para realizar estas tareas se estimaron 15 horas de trabajo pero en realidad fueron 18. Al finalizar el \eng{sprint} se completaron todas las tareas.

El \eng{burndown} de este \eng{sprint} se ilustra en la figura~\ref{fig:Burndown/Sprint3}.
\imagenflotante{Burndown/Sprint3}{\eng{Burndown chart} del \eng{sprint} 3.}

\subsection{\eng{Sprint} 4 (9/10/2025 -- 16/10/2025)}
Los objetivos del \myhref{https://github.com/lbr1014/TFG_RAG/milestone/4?closed=1}{Sprint 4} fueron continuar realizando la aplicación de \eng{Web Scraping}, tanto en \eng{Selenium} como en \eng{Playwright}, para sacar licitaciones de la página de \contsec{}. Incluir en la memoria una comparación entre ambas aplicaciones de \eng{Web Scraping}. Seguir investigando sobre los \eng{RAG}.
Para llevar a cabo estos objetivos se plantearon las siguientes tareas:
\begin{itemize}
\tightlist
	\item Continuar con el \eng{script} de \eng{web scraping} usando Playwright.
	\begin{itemize}
	\tightlist
		\item Continuar con el \eng{script} de \eng{web scraping} usando Selenium.
		\begin{itemize}
		\tightlist
			\item Estudiar sobre JSON.
		\end{itemize}
	\end{itemize}
	\item Escribir conceptos teóricos en la memoria.
	\begin{itemize}
	\tightlist
		\item Leer el capítulo 6 del libro: \eng{A simple guide to retrieval augmented generation}~\cite{book:RAG_SimpleGuide}.
	\end{itemize}
\end{itemize}

Para realizar estas tareas se estimaron 14 horas de trabajo pero en realidad fueron 18. Al finalizar el \eng{sprint} se completaron todas las tareas.

El \eng{burndown} de este \eng{sprint} se ilustra en la figura~\ref{fig:Burndown/Sprint4}.
\imagenflotante{Burndown/Sprint4}{\eng{Burndown chart} del \eng{sprint} 4.}

\subsection{\eng{Sprint} 5 (16/10/2025 -- 23/10/2025)}
Los objetivos del \myhref{https://github.com/lbr1014/TFG_RAG/milestone/5?closed=1}{Sprint 5} fueron investigar sobre las mejoras que aporta la \eng{API} asíncrona de \eng{Playwright} y como se implementa. También, se planificado investigar sobre la gestión de dependencias y entornos en \eng{Python}. Todo esto mientras se sigue investigando sobre los \eng{RAG} y \eng{LLM} y se escriben los aspectos más relevantes en la memoria. 
Para llevar a cabo estos objetivos se plantearon las siguientes tareas:
\begin{itemize}
\tightlist
	\item Realizar código con \eng{Playwright} asíncrono.
	\begin{itemize}
	\tightlist
		\item Investigar sobre la \eng{API} asíncrona de \eng{Playwright}.
	\end{itemize}
	\item Investigar sobre las dependencias y los entornos en \eng{Python}.
	\item Escribir conceptos teóricos en la memoria.
	\begin{itemize}
	\tightlist
		\item Leer más capítulo del libro: \eng{A simple guide to retrieval augmented generation}~\cite{book:RAG_SimpleGuide}.
	\end{itemize}
\end{itemize}

Para realizar estas tareas se estimaron 15 horas de trabajo pero en realidad fueron 18, ya que se tuvo que reescribir parte del código realizado en los \eng{sprints} anteriores para el \eng{Web Scraping} debido a una actualización en la página de \contsec{}. Este cambio afectó en mayor medida a la extracción de información de las licitaciones, ya que ha sido la página que más ha cambiado con dicha actualización. A pesar de este contratiempo al finalizar el \eng{sprint} se completaron todas las tareas.

El \eng{burndown} de este \eng{sprint} se ilustra en la figura~\ref{fig:Burndown/Sprint5}.
\imagenflotante{Burndown/Sprint5}{\eng{Burndown chart} del \eng{sprint} 5.}

\subsection{\eng{Sprint} 6 (23/10/2025 -- 30/10/2025)}

Los objetivos del \myhref{https://github.com/lbr1014/TFG_RAG/milestone/6?closed=1}{Sprint 6} fueron descargar los documentos (en formato PDF) de las licitaciones de la página de \contsec{}. Para descargar dichos PDF se va a usar de base el programa de \eng{Web Scraping} del \eng{sprint} anterior. Además, se han fijado como objetivos investigar sobre los modelos de inteligencia artificial de \eng{Hugging Face} e instalar y probar el funcionamiento de \eng{Ollama}.
Para llevar a cabo estos objetivos se plantearon las siguientes tareas:
\begin{itemize}
\tightlist
	\item Investigar sobre \eng{Ollama}.
	\item Investigar y probar modelos de inteligencia artificial de \eng{Hugging Face}.
	\item Investigar sobre formas de implementar los \eng{RAG}.
	\item Descargar los pdfs de las licitaciones de la la página de \contsec{}.
	\item Instalar las dependencias en \eng{Poetry}.
\end{itemize}

Para realizar estas tareas se estimaron 12 horas de trabajo pero en realidad fueron 15, se ha usado un servidor de la UBU para descargar un \ext{zip} con las licitaciones. Para acceder a dichas licitaciones se ha usado el archivo \texttt{<<resultados\_playwright\_asincrono>>} obtenido en el \eng{sprint} anterior, en el cual se almacenan los enlaces a los pdfs de cada licitación. Al finalizar el \eng{sprint} se completaron todas las tareas.

El \eng{burndown} de este \eng{sprint} se ilustra en la figura~\ref{fig:Burndown/Sprint6}.
\imagenflotante{Burndown/Sprint6}{\eng{Burndown chart} del \eng{sprint} 6.}

\subsection{\eng{Sprint} 7 (30/10/2025 -- 6/11/2025)}

El objetivo principal del \myhref{https://github.com/lbr1014/TFG_RAG/milestone/7?closed=1}{Sprint 7} fue investigar e implementar distintos tipos de \eng{tokenizer}. Esto se planteó como un primer paso para la implantación total del \eng{RAG}.
Para llevar a cabo estos objetivos se plantearon las siguientes tareas:
\begin{itemize}
\tightlist
	\item Investigar sobre los \eng{tokenizers}.
	\item Implementar distintos tipos de \eng{tokenizer}.
	\item Comparar las implementaciones de \eng{tokenizers}.
\end{itemize}

Para realizar estas tareas se estimaron 12 horas de trabajo pero en realidad fueron 15. En este \eng{sprint} se han llevado a cabo todas las tareas planificadas, pero sería destacable que en los \eng{scripts} de las diversas pruebas de \eng{tokenizers} no solo se han incluido estos. Si no que también se han incluido funciones para probar su uso, como pueden ser \eng{prompts} y funciones para dividir los pdfs de las licitaciones.

El \eng{burndown} de este \eng{sprint} se ilustra en la figura~\ref{fig:Burndown/Sprint7}.
\imagenflotante{Burndown/Sprint7}{\eng{Burndown chart} del \eng{sprint} 7.}

\subsection{\eng{Sprint} 8 (6/11/2025 -- 13/11/2025)}

El objetivo principal del \myhref{https://github.com/lbr1014/TFG_RAG/milestone/8?closed=1}{Sprint 8} fue implementar una base de datos para el \eng{RAG} y juntarlo con uno de los \eng{tokenizers} ya implementados del \eng{sprint} anterior. Además, en este \eng{sprint} se planteó como objetivo obtener los pliegos definitivos de de la página de . 

Esto se debe a que los pliegos obtenidos anteriormente no eran correctos, se trataban de resúmenes incompletos que no servían para alimentar al \eng{RAG}. Durante el \eng{sprint} 7, se encontraron los pliegos correctos y se fijo como objetivo del siguiente \eng{sprint} modificar los programa de \eng{Web Scraping} para descargarlos. 
Para llevar a cabo estos objetivos se plantearon las siguientes tareas:
\begin{itemize}
\tightlist
	\item Obtener pliegos correctos.
	\item Estudiar sobre las distintas bases de datos.
	\item Implementar la base de datos.
	\item Explicar la comparativa de las bases de datos.
\end{itemize}

Para realizar estas tareas se estimaron 12 horas de trabajo, pero en realidad fueron 16. En este \eng{sprint} se han llevado a cabo todas las tareas planificadas. Se ha implementado la base de datos vectorial \eng{Qdranty} y se ha juntado con el diseño de \eng{tokenizer} del \eng{sprint} anterior, aunque se ha cambiado el modelo para usar el del libro \myhref{https://www.amazon.com/LLM-Engineers-Handbook-engineering-production/dp/1836200072}{\eng{LLM Engineer's Handbook}}. Además, se han realizado los \eng{embeddings} para poder almacenarlos en la base de datos vectorial.

El \eng{burndown} de este \eng{sprint} se ilustra en la figura~\ref{fig:Burndown/Sprint8}.
\imagenflotante{Burndown/Sprint8}{\eng{Burndown chart} del \eng{sprint} 8.}

\subsection{\eng{Sprint} 9 (13/11/2025 -- 20/11/2025)}
El objetivo principal del \myhref{https://github.com/lbr1014/TFG_RAG/milestone/9?closed=1}{Sprint 9} fue realizar un \eng{script} que, dada una consulta, devuelva el \eng{chunk} más parecido. Junto con este \eng{chunk} debe devolver el nombre y el título del PDF, así como una respuesta del \eng{LLM} basándose en dicho \eng{chunk}. Para llevar a cabo estos objetivos se plantearon las siguientes tareas:
\begin{itemize}
\tightlist
	\item Implementar funciones que devuelvan el nombre del PDF al que pertenece el \eng{chunk}.
	\item Mejorar la partición de los \eng{chunks} añadiendo solapamiento (\eng{overlap}).
	\item Obtención del fragmento del \eng{chunk}.
	\item Implementar código para que el \eng{LLM} de una respuesta a partir de una pregunta del usuario.
	\item Añadir al código funciones para medir los tiempos.
	\item Usar \eng{loggers} en el código.
\end{itemize}

Para realizar estas tareas se estimaron 18 horas y más o menos se ha cumplido dicha estimación. En este \eng{sprint} se han llevado a cabo todas las tareas planificadas. Se ha mejorado el código añadiendo medidas de los tiempos de ejecución e implementando el uso de \eng{loggers}. También, se ha añadido solapamiento en los \eng{chunks} lo que mejora la eficacia de estos. Por otro lado, se ha implementado el objetivo principal, es decir, el \eng{script} que permite que el usuario realice una pregunta y el \eng{LLM} conteste basándose en el \eng{chunk} más parecido a la pregunta. Además se indica el pdf del que ha sacado la información asi como el \eng{chunk} concreto.

El \eng{burndown} de este \eng{sprint} se ilustra en la figura~\ref{fig:Burndown/Sprint9}.
\imagenflotante{Burndown/Sprint9}{\eng{Burndown chart} del \eng{sprint} 9.}

\subsection{\eng{Sprint} 10 (20/11/2025 -- 27/11/2025)}
El objetivo principal del \myhref{https://github.com/lbr1014/TFG_RAG/milestone/10?closed=1}{Sprint 10} fue intentar convertir los PDF a \eng{Markdown}. Durante la reunión de planificación del \eng{sprint} se decidió intentarlo para ver si el conversor a \eng{Markdown} es capaz de detectar la estructura jerárquica de los documentos PDF, ya que, si detectará las secciones y subsecciones se podrían añadir como metadatos. Esto mejoraría significativamente la relevancia de los \eng{chunks} al poder contextualizarlos mejor. Se han valorado diversas formas de realizar la conversión, por ejemplo, desde un \eng{OCR} vinculado a un \eng{LLM} o directamente empleando una biblioteca de \eng{Python}. Para llevar a cabo estos objetivos se plantearon las siguientes tareas:
\begin{itemize}
\tightlist
	\item Buscar formas de automatizar la conversión a \eng{Markdown}.
	\item Realizar la conversión de un pdf a \eng{Markdown}.
	\begin{itemize}
	\tightlist
		\item Automatizar la conversión a \eng{Markdown}.
	\end{itemize}
\end{itemize}

Para realizar estas tareas se estimaron 10 horas, pero debido a la gran heterogeneidad de formatos de los PDF fueron más, aproximadamente unas~15. A pesar de que se han obtenido resultados significativamente mejores que en las primeras aproximaciones, no se ha conseguido ningún \texttt{Markdown} lo suficientemente bueno. Esto se debe a que cada PDF de licitaciones tiene un formato distinto y no se ha conseguido ningún \eng{script} que detecte correctamente todas las secciones.

El \eng{burndown} de este \eng{sprint} se ilustra en la figura~\ref{fig:Burndown/Sprint10}.
\imagenflotante{Burndown/Sprint10}{\eng{Burndown chart} del \eng{sprint} 10.}

\subsection{\eng{Sprint} 11 (27/11/2025 -- 4/12/2025)}
El objetivo principal del \myhref{https://github.com/lbr1014/TFG_RAG/milestone/11?closed=1}{Sprint 11} fue intentar implementar un \eng{script} que basándose en los resultados obtenidos del anterior sea capaz de detectar correctamente las secciones de los PDF. Además, se propuso comenzar a investigar sobre el funcionamiento de \eng{Flask}. Ya que esta aplicación se va a utilizar en el proyecto para realizar una \eng{web} para realizar las consultas al \eng{LLM}. Para llevar a cabo estos objetivos se plantearon las siguientes tareas:
\begin{itemize}
\tightlist
	\item Probar nuevos OCR para la conversión de los PDF a \eng{Markdown}.
	\item Investigar sobre \eng{Flask}.
\end{itemize}

Para realizar estas tareas se estimaron 8 horas. Al final se emplearon cerca de 12 horas debido al mismo problema del \eng{sprint} anterior, que no se ha conseguido implementar un \eng{script} capaz de detectar correctamente todas las secciones.

El \eng{burndown} de este \eng{sprint} se ilustra en la figura~\ref{fig:Burndown/Sprint11}.
\imagenflotante{Burndown/Sprint11}{\eng{Burndown chart} del \eng{sprint} 11.}

\subsection{\eng{Sprint} 12 (4/12/2025 -- 15/01/2026)}
El objetivo principal del \myhref{https://github.com/lbr1014/TFG_RAG/milestone/12?closed=1}{Sprint 12} fue comenzar a realizar la página \eng{web} con \eng{Flask} y comenzar a definir los casos de uso del proyecto. Este \eng{sprint} fue más largo que el resto ya que coincidió con la semana de exámenes finales del primer cuatrimestre y las navidades. Para llevar a cabo estos objetivos se plantearon las siguientes tareas:
\begin{itemize}
\tightlist
	\item Definir los casos de uso.
	\begin{itemize}
	\tightlist
		\item Identificar los casos de uso del proyecto.
		\begin{itemize}
		\tightlist
			\item Identificar los requisitos del proyecto.
		\end{itemize}
	\end{itemize}
	\item Comenzar la página \eng{web}.
	\item Hacer la gestión básica de usuarios.
	\item Implementar la gestión de usuarios para los administradores.
	\item Hacer los \eng{test} para la gestión de usuarios de la \eng{web}.
\end{itemize}

Para realizar estas tareas se estimaron 28 horas de trabajo. Al final se emplearon más de 35 horas, debido principalmente a que la tarea <<Comenzar la página \eng{web}>> llevó más tiempo del planificado. Esta tarea consistió en realizar la estructura básica de la \eng{web}, la cual finalmente  se dividió en \eng{blueprints} para que tuviese mejor escalabilidad. Comprender el funcionamiento de los \eng{blueprints} y saber implementarlos aumentó significativamente el tiempo de desarrollo de la estructura de la \eng{web}.

El \eng{burndown} de este \eng{sprint} se ilustra en la figura~\ref{fig:Burndown/Sprint12}.
\imagenflotante{Burndown/Sprint12}{\eng{Burndown chart} del \eng{sprint} 12.}

\subsection{\eng{Sprint} 13 (15/01/2025 -- 22/01/2026)}
El objetivo principal del \myhref{https://github.com/lbr1014/TFG_RAG/milestone/13?closed=1}{Sprint 13} fue completar las funcionalidades de la \eng{web} permitiendo que los administradores puedan cargar documentos, de forma manual o automática (mediante \eng{web scraping}). El sistema gestionará estos documentos para extraer de ellos los \eng{embeddings} necesarios para el \eng{RAG}. Además, se marcó como objetivo incluir la interfaz para realizar las consultas al \eng{LLM}. Para llevar a cabo estos objetivos se plantearon las siguientes tareas:
\begin{itemize}
\tightlist
	\item Añadir casos de uso.
	\item Incluir una clase consulta.
	\item Añadir la funcionalidad carga de documentos.
	\item Mostrar historial de consultas.
	\item Implementar las consultas en la \eng{web Flask}.
	\item Juntar el \eng{web scraping} con la carga de documentos de la \eng{web}.
	\item Realizar la interfaz de las consultas.
	\item Poner la clave secreta como variable de entorno.
\end{itemize}

Para realizar estas tareas se estimaron 28,5 horas de trabajo. Al final del \eng{sprint} se cumplieron todas las tareas más o menos en el plazo esperado.

El \eng{burndown} de este \eng{sprint} se ilustra en la figura~\ref{fig:Burndown/Sprint13}.
\imagenflotante{Burndown/Sprint13}{\eng{Burndown chart} del \eng{sprint} 13.}

\subsection{\eng{Sprint} 14 (22/01/2025 -- 29/01/2026)}

El objetivo principal del \myhref{https://github.com/lbr1014/TFG_RAG/milestone/14?closed=1}{Sprint 14} fue convertir los procesos de \eng{web scraping} y de actualización de la base de datos vectorial a modo asíncrono. En la reunión se decidió priorizar este cambio, ya que la implementación anterior, en modo síncrono, tardaba mucho tiempo y bloqueaba la web para el usuario. Por otro lado, se decidió continuar desarrollando la documentación, concretamente los anexos. Para llevar a cabo estos objetivos se plantearon las siguientes tareas:
\begin{itemize}
\tightlist
	\item Cambios en las clases de la aplicación.
	\begin{itemize}
	\tightlist
		\item Cambios en la clase documentos.
	\end{itemize} 
	\item Escribir en los anexos el diseño de datos.
	\begin{itemize}
	\tightlist
		\item Realizar el diagrama relacional.
		\item Realizar el diagrama Entidad/Relación.
	\end{itemize}
	\item Aplicar correcciones a los anexos.
	\item Realizar el diagrama de casos de uso.
	\item Diseño de interfaces.
	\item Cambio de funcionalidades a modo asíncrono.
	\item Incluir mensajes de \eng{feedback} para el usuario.
	\item Incluir las clases \code{Chunk} y \code{Embedding}.
	\item Aplicar cambios en el código para integrar correctamente las nuevas clases.
	\begin{itemize}
	\tightlist
		\item Crear tabla \code{ConsultaChunk}.
	\end{itemize}
	\item Revisar y corregir código usando \eng{SonarCloud}.
	\item Hacer que el \eng{LLM} use más de un \eng{chunk} para responder.
\end{itemize}

Para realizar estas tareas se estimaron 50 horas de trabajo. Al final del \eng{sprint} se cumplieron todas las tareas. El tiempo final dedicado fue más o menos el estimado. A pesar de que algunas tareas (como las de corregir el código usando \eng{SonarCloud} y las de cambiar las funcionalidades a modo asíncrono) se retrasaron, hubo otras como las de realizar los diagramas que se adelantaron.

El \eng{burndown} de este \eng{sprint} se ilustra en la figura~\ref{fig:Burndown/Sprint14}.
\imagenflotante{Burndown/Sprint14}{\eng{Burndown chart} del \eng{sprint} 14.}

\subsection{\eng{Sprint} 15 (29/01/2025 -- 05/02/2026)}
El objetivo principal del \myhref{https://github.com/lbr1014/TFG_RAG/milestone/15?closed=1}{Sprint 15} fue levantar mi aplicación \eng{Flask} usando un contenedor \eng{docker}. En la reunión se decidió implementar un \eng{docker} ya que mejora la potabilidad, la ligereza y el aislamiento de la aplicación. Para llevar a cabo estos objetivos se plantearon las siguientes tareas:
\begin{itemize}
\tightlist
	\item Implementar un \eng{docker} para la aplicación \eng{Flask}.
	\begin{itemize}
	\tightlist
		\item Estudiar como desplegar la aplicación \eng{web}.
		\item Arreglar la indexación de documentos para usar el \eng{docker}.
		\item Arreglo de \eng{Web Scraping} para evitar fallos por perdida de conexión.
	\end{itemize}
	\item Empezar a escribir la viabilidad económica.
	\item Completar los \eng{test} unitarios realizados previamente.
	\item Realiza \eng{test} unitarios para las nuevas clases.
\end{itemize}

Para realizar el \eng{sprint} se estimaron 34 horas de trabajo. Al final de este \eng{sprint} se invirtieron unas 40 horas, pero, no se completaron todas las tareas, quedando incompletas las tareas <<Completar los \eng{test} unitarios realizados previamente>> y <<Realiza \eng{test} unitarios para las nuevas clases>>. Esto se debió a que la implementación de \eng{docker} llevo más tiempo del esperado por fallos en la conexión con el módulo de \eng{qdrant}. Las tareas incompletas se van a llevar a cabo en el siguiente \eng{sprint}.

El \eng{burndown} de este \eng{sprint} se ilustra en la figura~\ref{fig:Burndown/Sprint15}.
\imagenflotante{Burndown/Sprint15}{\eng{Burndown chart} del \eng{sprint} 15.}

\subsection{\eng{Sprint} 16 (05/02/2025 -- 12/02/2026)}
El objetivo principal del \myhref{https://github.com/lbr1014/TFG_RAG/milestone/16?closed=1}{Sprint 16} fue incluir el uso de \eng{Nginx} para levantar la aplicación \eng{web} en el servidor integrándolo con la implementación de \eng{Gunicode} del \eng{sprint} anterior. Además, se fijo como objetivo explicar en la memoria algunos trabajos relacionados. Para llevar a cabo estos objetivos se plantearon las siguientes tareas:
\begin{itemize}
\tightlist
	\item Incluir trabajos relacionados en la memoria.
	\item Incluir imágenes en la memoria.
	\item Correcciones documentación.
	\item Completar algunos de los \eng{tests} unitarios realizados previamente.
	\item Realizar \eng{tests} unitarios para las nuevas clases.
	\item Mejoras en el \eng{docker}.
	\item Incluir \eng{Nginx} en la aplicación \eng{Flask}.
\end{itemize}

Para realizar el \eng{sprint} se estimaron 34 horas de trabajo. Al finalizar el \eng{sprint} se realizaron todas las tareas, y de manera general respetando la planificación temporal. En el caso de la tarea <<Incluir trabajos relacionados en la memoria>> se estimó que se iba a realizar en 4 horas pero finalmente fueron unas 6 horas, ya que, se realizó una comparativa entre ellos y con el proyecto que se esta llevando a cabo en este TFG.

El \eng{burndown} de este \eng{sprint} se ilustra en la figura~\ref{fig:Burndown/Sprint16}.
\imagenflotante{Burndown/Sprint16}{\eng{Burndown chart} del \eng{sprint} 16.}

\subsection{\eng{Sprint} 17 (12/02/2025 -- 19/02/2026)}
El objetivo principal del \myhref{https://github.com/lbr1014/TFG_RAG/milestone/17?closed=1}{Sprint 17} fue incluir que la aplicación pueda mandar correos a los usuarios, concretamente para avisar cuando acaban de ejecutarse los procesos largos (como el \eng{web scraping} o la actualización de la base de datos vectorial) y para la recuperación de la contraseña cuando se le ha olvidado. También, se siguieron implementando los \eng{tests} unitarios de la aplicación para cubrir el \eng{coverage}. Para llevar a cabo estos objetivos se plantearon las siguientes tareas:
\begin{itemize}
\tightlist
	\item Correo cuando finalicen las tareas largas.
	\begin{itemize}
	\tightlist
		\item Incluir recuperación de contraseña.
		\begin{itemize}
		\tightlist
			\item Investigar como se pueden mandar correos automáticos desde \eng{Flask}.
		\end{itemize}
	\end{itemize}
	\item Completar \eng{tests} y generar el informe del \eng{coverage}.
	\item Incluir conversión a \eng{markdown} en la aplicación.
	\item Botón cancelar consulta.
\end{itemize}

Para realizar el \eng{sprint} se estimaron 24 horas de trabajo. Al finalizar el \eng{sprint} se realizaron todas las tareas invirtiendo unas 30 horas. Esto se debió a que la realización de los \eng{tests} llevó más tiempo del planificado. Sobre todo, aquellos los de las clases concurrentes y aquellos para los que se necesitaba entorno con bases de datos.

El \eng{burndown} de este \eng{sprint} se ilustra en la figura~\ref{fig:Burndown/Sprint17}.
\imagenflotante{Burndown/Sprint17}{\eng{Burndown chart} del \eng{sprint} 17.}

\subsection{\eng{Sprint} 18 (19/02/2025 -- 26/02/2026)}


\subsection{\eng{Sprint} 19 (26/02/2025 -- 5/03/2026)}

\section{Estudio de viabilidad}
El estudio de viabilidad es un análisis detallado que se realiza antes de iniciar un proyecto. Sirve para determinar la rentabilidad de ese proyecto, es decir, si se pude hacer realidad. Para ello se analiza si es económicamente rentable (sección~\hyperref[viabilidad_economica]{A.3.1}) y legalmente viable (sección~\hyperref[viabilidad_legal]{A.3.2}). 

\subsection{Viabilidad económica}\label{viabilidad_economica}
Para analizar la viabilidad económica se ha seguido el coste total de propiedad (TCO)~\cite{book:TCO_costs}. El TCO consiste en identificar y estimar todos los costes asociados al ciclo de vida de un producto tecnológico, no solo los costes directos de su adquisición, sino también los derivados de su operación, mantenimiento y uso a lo largo del tiempo. De esta forma, se obtiene una visión global y realista del impacto económico del sistema. Los costes considerados vienen recogidos en la tabla~\ref{tabla:explicacion_TCO}.

\begin{table}[h]
	\centering
	\begin{tabularx}{\linewidth}{
    		>{\raggedright\arraybackslash}p{0.46\linewidth}
    		>{\raggedleft\arraybackslash}p{0.46\linewidth}}
		\toprule
		\textbf{Componentes}    & \textbf{Componentes del costo}\\
		\midrule
		Adquisición \eng{hardware} & Precio de compra del equipo \eng{hardware}. Se debe tener en cuenta los ordenadores, terminales, almacenamiento e impresoras. \\
		Adquisición \eng{software} & Compra de licencias y \eng{software} para cada usuario.  \\
		Instalación & Coste para instalar los ordenadores y el \eng{software}. \\
		Capacitación & Coste de entrenamiento de los especialistas y de los usuarios finales. \\
		Soporte & Coste para dar soporte técnico continuo. \\
		Mantenimiento & Coste para mantener y actualizar el \eng{hardware} y \eng{software}. \\
		Infraestructura & Coste por adquirir, mantener y dar soporte a la infraestructura relacionada. Incluye las redes y el equipo especializado. \\
		Tiempo inactivo & Coste de pérdidas de productividad debido a fallos de \eng{hardware} o \eng{software}. Estos fallos provocan que el sistema no este disponible. \\
		Espacio y energía & Coste de bienes y servicios públicos para aojar y proveer energía para la tecnología. \\
		\bottomrule
	\end{tabularx}
	\caption{Coste total de desarrollo (TCO).}
	\label{tabla:explicacion_TCO}
\end{table}

Para llevar a cabo el estudio de la viabilidad económica se va a considerar que el trabajo llevado a cabo por la alumna lo ha realizado un desarrollador contratado durante los 9 meses que ha durado el proyecto. Si se supone un escenario empresarial real en España y se toma como referencia el \myhref{https://www.boe.es/buscar/doc.php?id=BOE-A-2024-2251}{Salario Mínimo Interprofesional (SMI) de 2024}, el salario mínimo anual es de 15\,876\,€, es decir, 1\,134\,€ al mes (se divide en 14 pagas). Además, a este salario hay que añadirle los gastos por la cotización a la \myhref{https://www.seg-social.es/wps/portal/wss/internet/Trabajadores/CotizacionRecaudacionTrabajadores/36537}{Seguridad Social}. Normalmente la cotización es del 28,30\%, de los cuales 23,60\% se corresponden con gastos de la empresa. Por lo cual el gasto por la contratación del desarrollado sería:
$$
1\,134\frac{\,\text{\texteuro}}{\,\text{meses}} \cdot 1,236 \cdot 9\,\text{meses}
= 12\,614,62\,\text{\texteuro}
$$

Respecto a la adquisición de \eng{hardware} no ha sido necesaria la adquisición de \eng{hardware} específico adicional. El sistema \eng{RAG} se ha implementado de forma que se pueda ejecutar en entornos estándares, sin requerir servidores dedicados ni equipamiento especializado. El desarrollo y las pruebas se han realizado utilizando mi equipo personal y servidores ya disponibles en la universidad. El equipo utilizado ha sido un \myhref{https://support.hp.com/es-es/document/ish_7878690-7886844-16}{\eng{Victus by HP Gaming Laptop 16-r0xxx}}. El coste de adquisición del portátil fueron unos 1\,200\,€. Considerando un periodo de amortización de 5 años, el coste amortizado correspondiente al proyecto con una duración de 9 meses (del 11 de septiembre al 11 junio) sería:
$$
\left(\frac{1\,200\,\text{\texteuro}}{5\,\text{años} \cdot 12\,\text{meses}}\right)\cdot 9\,\text{meses}
= 180\,\text{\texteuro}
$$

Además, se ha usado el servidor \myhref{https://www.gigabyte.com/es/Enterprise/Server-Motherboard/MW51-HP0-rev-1x}{\eng{beta}} de la universidad. Aunque el servidor no se ha considerado un gasto de \eng{hardware} ya que se ha considerado que se amortizaba en 5 años y ya tienen 9.

Respecto a la adquisición de \eng{software} la mayor parte del \eng{software} utilizado en el proyecto es de código abierto y de uso gratuito, como \eng{Python, Flask, Qdrant, PostgresSQL, Docker} o las librerías de \eng{Hugging Face}. Estas herramientas no implican un coste directo de licencia. Por lo cual el único gasto que debe considerarse en este apartado es la adquisición de la licencia de 
\myhref{https://www.microsoft.com/es-es/d/windows-11-home/dg7gmgf0krt0/000P}{\eng{Windows}} con un coste de unos 145\,€. Se considera una amortización de 5 años, por lo cual para los 9 meses del proyecto:
$$
\left(\frac{145\,\text{\texteuro}}{5\,\text{años} \cdot 12\,\text{meses}}\right)\cdot 9\,\text{meses}
= 21,75\,\text{\texteuro}
$$
Por otro lado en el proyecto no ha habido costes indirectos asociados a servicios de terceros como pueden ser las API o modelos en la nube, ya que se han elegido alternativas gratuitas.

Los costes de instalación incluyen el tiempo invertido en la configuración del entorno de desarrollo, instalación de dependencias, despliegue del sistema con \eng{Docker} y configuración de las bases de datos. Este tiempo debe considerar como un coste laboral. Se estima que en estas tareas se ha invertido un 20\% del tiempo total del proyecto. Por lo cual teniendo en cuenta que el coste laboral del desarrollador ha sido de 12\,614,62\,€ el de los costes de instalación se corresponde:
$$
12\,614,62\,\text{\texteuro} \cdot 0,20
= 2\,522,92\,\text{\texteuro}
$$

Respecto a los costes de capacitación, se ha requerido una fase importante de aprendizaje en tecnologías como \eng{RAG}, bases de datos vectoriales, modelos \eng{LLM}, \eng{web scraping} y despliegue \eng{web}. Este tiempo de formación constituye un coste indirecto de capital humano. Se estima que aproximadamente un 25\% del tiempo total del proyecto ha sido dedicado a tareas de aprendizaje. Por lo cual, el coste de capacitación ha sido:
$$
12\,614,62\,\text{\texteuro} \cdot 0,25
= 3\,153,66\,\text{\texteuro}
$$

El soporte durante la ejecución del proyecto ha sido llevado a cabo por los tutores. Se han llevado a cabo reuniones (\eng{sprints}) para la revisión de decisiones y la resolución de problemas. Han sido necesarias unas 28 reuniones a lo largo del proyecto de aproximadamente una hora cada una. Además, al tiempo invertido para el soporte técnico, debe sumarse el tiempo para llevar a cabo las revisiones de la memoria, los anexos y del código. En total se estima que los tutores han invertido unas 100 horas en el proyecto cada uno. Se va a suponer que su sueldo es de 20\,€ la hora al mes sería:
$$ 
20\frac{\,\text{\texteuro}}{\,\text{hora}} \cdot 7\frac{\,\text{horas}}{\,\text{días}} \cdot 5\frac{\,\text{días}}{\,\text{semana}} \cdot 4\frac{\,\text{semanas}}{\,\text{mes}} 
= 2\,800\frac{\,\text{\texteuro}}{\,\text{mes}}
$$
Por lo cual el gasto en soporte sería:
$$ 
20\frac{\,\text{\texteuro}}{\,\text{hora}} \cdot 200\,\text{horas}
= 4\,000\,\text{\texteuro}
$$

En el mantenimiento se incluye corrección de errores, adaptación a cambios en las fuentes de obtención de datos (como las modificaciones en la \eng{web} de contratación pública) y actualización de dependencias. A lo largo del proyecto se han producido varias modificaciones del código motivadas por cambios externos, por lo cual ha sido necesario dedicar tiempo periódico al mantenimiento correctivo y evolutivo. Este coste se traduce en un 5\% del tiempo total de desarrollo del proyecto. Por lo cual el coste de mantenimiento ha sido:
$$
12\,614,62\,\text{\texteuro} \cdot 0,05
= 630,73\,\text{\texteuro}
$$

Respecto a la infraestructura necesaria para el proyecto, se basa en servicios de red, bases de datos y servidores locales. Se han utilizado infraestructuras universitarias (el servidor cuya amortización ya se ha tenido en cuenta en el apartado de adquisición de \eng{hardware}) y servicios gratuitos, por lo cual no es necesario sumar gastos e infraestructuras al proyecto. 

El tiempo inactivo se corresponde al impacto económico que suponen los periodos en los que el sistema no está operativo o no puede utilizarse de forma eficiente.

Finalmente, deben considerarse los costes asociados al espacio físico del trabajo y al consumo de energía del equipo. Se ha considerado que para llevar a cabo el proyecto se ha gastado una media de 80\,€ mensuales de electricidad, como el proyecto ha durado 9 meses, el coste total sería:
$$
80\frac{\,\text{\texteuro}}{\,\text{meses}} \cdot 9\,\text{meses}
= 720\,\text{\texteuro}
$$

Los costes totales del proyecto vienen desglosados en la siguiente tabla~\ref{tabla:coste}

\begin{table}[h]
	\centering
	\begin{tabularx}{\linewidth}{
    		>{\raggedright\arraybackslash}p{0.46\linewidth}
    		>{\raggedleft\arraybackslash}p{0.46\linewidth}}
		\toprule
		\textbf{Componentes}    & \textbf{Estimación (euros)}\\
		\midrule
		Adquisición \eng{hardware} & 180  \\
		Adquisición \eng{software} & 21,75  \\
		Instalación & 2\,522,92  \\
		Capacitación & 3\,153,66 \\
		Soporte &  4\,000 \\
		Mantenimiento & 630,73 \\
		Infraestructura & 0 \\
		Tiempo inactivo & 0 \\
		Espacio y energía & 720 \\
		Coste total del proyecto & 11\,229 \\
		\bottomrule
	\end{tabularx}
	\caption{Coste total del proyecto.}
	\label{tabla:coste}
\end{table}

\subsection{Viabilidad legal}\label{viabilidad_legal}


